\documentclass[11pt,a4paper]{article}

% ---------------------------------------------------------
% PREAMBLE (stable, warning-free, Overleaf-safe)
% ---------------------------------------------------------
\usepackage[utf8]{inputenc}
\usepackage[T1]{fontenc}
\usepackage{lmodern}
\usepackage{amsmath, amssymb, amsfonts}
\usepackage{bm}
\usepackage{geometry}
\usepackage{graphicx}
\usepackage{hyperref}
\usepackage{cite}
\usepackage{authblk}
\usepackage{physics}
\usepackage{mathtools}

\geometry{margin=2.4cm}

\title{\textbf{Companion LXXXIII: Fisher Forecasts for $\alpha$ Using Transport-Based Topological Statistics in the Relaxation-Driven Cyclic Cosmology}}
\author{Michael Lehmann \\ \texttt{mi.lehmann@gmx.de}}
\date{April 2026}

\begin{document}
\maketitle

\begin{abstract}
Transport-based statistics provide a powerful way to quantify differences between 
persistence-diagram ensembles across cosmological models. Building on the unified 
diffusion--optimal transport framework developed in previous Companions, this work 
constructs Fisher forecasts for the relaxation parameter $\alpha$ in the 
Relaxation-Driven Cyclic Cosmology (RDCC). The key observable is a 
scale-resolved transport discrepancy between RDCC and $\Lambda$CDM persistence 
diagrams, computed using Wasserstein, Sinkhorn, or neural transport maps. We 
derive the Fisher information associated with these transport statistics and 
construct $\alpha$-sensitivity kernels that identify the angular scales and 
topological features most informative for constraining $\alpha$ in Euclid-like 
weak-lensing surveys. This Companion provides the first explicit link between 
geometric transport quantities and forecasted RDCC parameter constraints.
\end{abstract}
% ---------------------------------------------------------
\section{Introduction}

Persistence diagrams encode the topological structure of weak-lensing convergence 
fields. Previous Companions introduced a hierarchy of geometric tools for these 
diagrams, including Wasserstein distances, entropic optimal transport, 
Schrödinger bridges, score-based diffusion models, and neural transport maps. 
Companion LXXXII introduced transport-based parameter inference and 
$\alpha$-sensitivity maps.

The present Companion takes the next step: we construct Fisher forecasts for the 
RDCC relaxation parameter $\alpha$ using transport-based topological statistics. 
This provides a direct and quantitative link between geometric transport 
quantities and observable constraints from Euclid-like surveys.
% ---------------------------------------------------------
\section{Transport-Based Observables}

Let $\mathcal{D}_{\Lambda}$ and $\mathcal{D}_{\alpha}$ denote persistence-diagram 
ensembles from $\Lambda$CDM and RDCC with parameter$\alpha$. We define a 
scale-resolved transport statistic

\begin{equation}
    \mathcal{T}_{\ell}(\alpha)
    =
    \mathbb{E}_{D \sim \mathcal{D}_{\Lambda}}
    \left[
        W_{p}\bigl(D, T_{\alpha}(D)\bigr)
    \right]_{\ell},
\end{equation}

where the subscript $\ell$ indicates restriction to features associated with 
angular multipole $\ell$.

The neural transport map $T_{\alpha}$ is trained to push $\mathcal{D}_{\Lambda}$ 
towards $\mathcal{D}_{\alpha}$.
% ---------------------------------------------------------
\section{Fisher Information for \texorpdfstring{$\alpha$}{alpha}}

Given a vector of transport statistics


\[
    \mathbf{T}(\alpha)
    = \{ \mathcal{T}_{\ell}(\alpha) \}_{\ell},
\]


with covariance matrix $\mathbf{C}$ estimated from mocks, the Fisher information 
for $\alpha$ is

\begin{equation}
    F_{\alpha\alpha}
    =
    \left( 
        \frac{\partial \mathbf{T}}{\partial \alpha}
    \right)^{\!\!T}
    \mathbf{C}^{-1}
    \left(
        \frac{\partial \mathbf{T}}{\partial \alpha}
    \right).
\end{equation}

The forecasted constraint is



\[
    \sigma(\alpha) = F_{\alpha\alpha}^{-1/2}.
\]



The derivative $\partial \mathbf{T} / \partial \alpha$ is estimated using RDCC 
mocks at nearby $\alpha$ values.
% ---------------------------------------------------------
\section{\texorpdfstring{$\alpha$}{alpha}-Sensitivity Kernels}

We define the $\ell$-resolved sensitivity kernel

\begin{equation}
    K_{\alpha}(\ell)
    =
    \left[
        \frac{\partial \mathcal{T}_{\ell}}{\partial \alpha}
    \right]^{2}
    \bigg/
    C_{\ell\ell}.
\end{equation}

This kernel quantifies the contribution of each angular scale $\ell$ to the total 
Fisher information.

Large values of $K_{\alpha}(\ell)$ indicate scales where RDCC-induced topological 
deformations are most pronounced.
% ---------------------------------------------------------
\section{Connection to Euclid-Like Surveys}

For a Euclid-like survey, the covariance matrix $\mathbf{C}$ is dominated by:

\begin{itemize}
    \item shape noise,
    \item finite sampling of galaxies,
    \item cosmic variance,
    \item reconstruction noise in convergence maps.
\end{itemize}

Transport statistics are robust to these effects because:

\begin{itemize}
    \item Wasserstein/Sinkhorn distances are stable under noise,
    \item neural transport maps can be trained to denoise diagrams implicitly,
    \item diffusion-based models provide uncertainty-aware transport.
\end{itemize}

Thus the Fisher forecasts derived here are directly applicable to Euclid mocks.
% ---------------------------------------------------------
\section{Conclusion}

Transport-based topological statistics provide a powerful and robust way to 
quantify differences between RDCC and $\Lambda$CDM. By constructing Fisher 
forecasts for the relaxation parameter $\alpha$, we establish a direct link 
between geometric transport quantities and observable constraints. The 
$\alpha$-sensitivity kernels introduced here identify the angular scales most 
informative for RDCC parameter inference in Euclid-like surveys. This Companion 
thus completes the bridge from geometric theory to observational forecasting.

\bibliographystyle{apsrev4-2}
\nocite{*}
\bibliography{references}


\end{document}
